%--------------------------------------
%CANEVAS
%--------------------------------------

%utiliser les environnement \begin{comment} \end{comment} pour mettre en commentaire le préambule une fois la programmation appelée dans le document maître (!ne pas oublier de mettre en commentaire \end{document}!)

\begin{comment}

\documentclass[a4paper, 11pt, twoside, fleqn]{memoir}

\usepackage{AOCDTF}

%--------------------------------------
%CANEVAS
%--------------------------------------

\newcommand\BoxColor{\ifcase\thechapshift blue!30\or brown!30\or pink!30\or cyan!30\or green!30\or teal!30\or purple!30\or red!30\or olive!30\or orange!30\or lime!30\or gray!\or magenta!30\else yellow!30\fi} %définition de la couleur des marqueurs de chapitre

\newcounter{chapshift} %compteur de chapitre du marqueur de chapitre
\addtocounter{chapshift}{-1}
	
\newif\ifFrame %instruction conditionnelle pour les couleurs des pages
\Frametrue

\pagestyle{plain}

% the main command; the mandatory argument sets the color of the vertical box
\newcommand\ChapFrame{%
\AddEverypageHook{%
\ifFrame
\ifthenelse{\isodd{\value{page}}}
  {\backgroundsetup{contents={%
  \begin{tikzpicture}[overlay,remember picture]
  \node[
  	rounded corners=3pt,
    fill=\BoxColor,
    inner sep=0pt,
    rectangle,
    text width=1.3cm,
    text height=5.5cm,
    align=center,
    anchor=north west
  ] 
  at ($ (current page.north west) + (-0cm,-2*\thechapshift cm) $) %nombre négatif = espacement des marqueurs entre les différents chapitres (à régler en fin de rédaction) (4.5cm vaut un espacement équivalement à la hauteur du marqueur, une page peut en contenir 6 avec cet espacement-la mais il est le plus équilibré)
    {\rotatebox{90}{\hspace*{.3cm}%
      \parbox[c][0.9cm][t]{5cm}{%
        \raggedright\textcolor{black}{\sffamily\textbf{\leftmark}}}}};
  \end{tikzpicture}}}
  }
  {\backgroundsetup{contents={%
  \begin{tikzpicture}[overlay,remember picture]
  \node[
  	rounded corners=3pt,
    fill=\BoxColor,
    inner sep=0pt,
    rectangle,
    text width=1.3cm,
    text height=5.5cm,
    align=center,
    anchor=north east
  ] 
  at ($ (current page.north east) + (-0cm,-2*\thechapshift cm) $) %nombre négatif = espacement des marqueurs entre les différents chapitres (à régler en fin de rédaction) (4.5cm vaut un espacement équivalement à la hauteur du marqueur, une page peut en contenir 6 avec cet espacement-la mais il est le plus équilibré)
    {\rotatebox{90}{\hspace*{.3cm}%
      \parbox[c][0.9cm][t]{5cm}{%
        \raggedright\textcolor{black}{\sffamily\textbf{\leftmark}}}}};
  \end{tikzpicture}}}%
  }
  \BgMaterial%
  \fi%
}%
  \stepcounter{chapshift}
}

\renewcommand\chaptermark[1]{\markboth{\thechapter.~#1}{}} %redéfinition du marqueur de chapitre pour ne contenir que le titre du chapitre %à personnaliser selon le nombre de chapitre dans le cours

%--------------------------------------
%corps du document
%--------------------------------------

\begin{document} %corps du document
	\openleft %début de chapitre à gauche

\end{comment}

\begingroup

%--------------------------------------
%logo Compagnon
%--------------------------------------
 
\begin{tikzpicture}[remember picture, overlay]
	\begin{scope}[shift={(current page.north west)}]
		\draw (1,-1)
		node [below right]{\includegraphics[scale=0.1]{logo_compagnons_nom.png}};
	\end{scope}
\end{tikzpicture}

~\\[5cm] %espace insécable pour marquer le début du texte sur l'environnement et situer d'autres éléments de texte sur la page

%--------------------------------------
%titre de la page
%--------------------------------------

\begin{flushleft}
	\HUGE\sffamily\textbf{Matière}\\
\end{flushleft}
\HRule
\begin{flushright}
	\huge\sffamily\textbf{Cours}\\[0.4cm]
\end{flushright}
~\\[2cm]

%--------------------------------------
%auteur et logo MTA
%--------------------------------------

\begin{minipage}{0.3\textwidth}
 \noindent\large
		\rmfamily Prénom1 \textsc{Nom1}
		\\
		\rmfamily Prénom2 \textsc{Nom2}
\end{minipage}
\hfill
\begin{minipage}{0.3\textwidth}
	\centering
	\includegraphics[scale=0.25]{logo_metier.png}
\end{minipage}

\begin{comment}
\hfill
\begin{minipage}{0.3\textwidth}
	\begin{flushright} \large
		\rmfamily Prénom3 \textsc{Nom3} %auteur supplémentaires à rajouter en mettant en commentaire l'environnement comment
		\\
 		\rmfamily Prénom4 \textsc{Nom4}
	\end{flushright}
\end{minipage}
\end{comment}

\vfill
\begin{flushleft}
	\sffamily\bfseries\large\ Édition \the\year.\the\month
\end{flushleft}
\vfill

%--------------------------------------
%pied de page Compagnon
%--------------------------------------

\begin{tikzpicture}[remember picture, overlay]
	\begin{scope}[shift={(current page.south west)}]
		\draw (1,0)
		node [above right]{\includegraphics{pied_page_compagnons.png}};
	\end{scope}
\end{tikzpicture}

\endgroup
\null\clearpage

%\end{document}